\chapter{Математика retrieval: от локального шага к глобальному поведению}
\label{ch:retrievalmath}

\section{Зачем отдельная математическая глава}

В базовом изложении retrieval уже описан корректно: есть якоря, есть обход,
есть настроение, есть телепорты, есть разные политики. Но для курса этого мало.
Нужно увидеть разницу между \emph{локальным законом шага} и \emph{глобальной
физикой потока}. Именно здесь многие архитектурные интуиции либо укрепляются,
либо рассыпаются.

\section{Переходный оператор}

Пусть граф памяти задан вершинами $V$ и типизированными рёбрами $E$. Для фиксированного
режима работы зададим для пары соседних узлов $A,B$ ненормированную силу перехода
\[
S(A\to B)
=
w_{AB}\cdot
\mathrm{rel}(B\mid q)\cdot
\mathrm{mc}(B\mid \mood)\cdot
\mathrm{unc}(B)^{\beta_u}\cdot
\mathrm{gate}(A,B),
\]
где последний множитель $\mathrm{gate}(A,B)$ удобно выделить отдельно: он собирает
всё, что не является содержательной близостью, но влияет на сам факт допустимости
шага — например, запрет на некоторые типы рёбер в целевом retrieval, штраф за
повторное хождение по одному и тому же контуру, локальные бюджеты.

Нормировка даёт матрицу переходов
\[
P_{AB}=\frac{S(A\to B)}{\sum\limits_{C\in \Neigh(A)} S(A\to C)}.
\]
Если включён телепорт, фактический оператор имеет вид
\[
\widetilde P=(1-\alpha)P+\alpha \mathbf{1}r^\top,
\]
где $r$ — распределение рестарта. В retrieval оно сосредоточено на якорях, во
внутренней жизни смещено к зонам неопределённости, во сне может быть почти
равномерным или управляться картой неудовлетворённого покрытия.

\begin{intuition}[Почему важен именно оператор, а не набор красивых факторов]
Пока формула живёт как список слов — релевантность, настроение, новизна, — легко
вообразить, что эти ручки независимы. Но после нормировки они начинают бороться
за один и тот же переходный вес. Усилив настроение, можно случайно ухудшить
релевантность; усилив жадность, можно обрушить покрытие. У оператора есть
системные эффекты, которых не видно по отдельным множителям.
\end{intuition}

\section{Температура и режим выбора}

Пусть логарифм ненормированного веса шага равен $z_{AB}=\log S(A\to B)$. Тогда
softmax с температурой $T$ задаёт
\[
P_{AB}(T)=\frac{\exp(z_{AB}/T)}{\sum_C \exp(z_{AC}/T)}.
\]
Пределы понятны:
\[
\begin{aligned}
T &\to 0      && \Rightarrow \text{почти жадный выбор},\\
T &\to \infty && \Rightarrow \text{почти равномерный выбор по допустимым соседям}.
\end{aligned}
\]

Но practically важен не предел, а производная: насколько чувствительна вероятность
перехода к малому изменению веса. Именно здесь видно, что высокая температура — не
просто «больше случайности», а меньшая устойчивость к мелким локальным ошибкам в
оценках. Для debugging это полезно, для production-retrieval — опасно.

\section{Якоря, рестарты и two-step retrieval}
\label{sec:retrievalmathtax}

В исходной архитектуре уже сказано, что retrieval не ищет запись обходчиком с нуля,
а сначала находит якоря. Это место стоит формализовать.

Пусть есть запрос $q$. Тогда retrieval распадается на две задачи:
\[
q \longmapsto A(q)\subseteq V
\longmapsto
\text{локальное разворачивание окрестности}.
\]
Первая задача — \emph{anchor selection}; вторая — \emph{context expansion}. Смешивать
их нельзя.

\begin{definition}[Anchor recall]
Anchor recall — доля запросов, для которых на этапе выбора якорей в стартовое
множество попал хотя бы один узел из минимального носителя правильного ответа.
\end{definition}

\begin{definition}[Bridge rate]
Bridge rate — доля разворачиваний, в которых после попадания в правильную область
обход действительно дошёл до промежуточных узлов, без которых ответный фрагмент
остался бы несвязной кучей фактов.
\end{definition}

Эти две метрики принципиально различны. Низкий anchor recall чинится поиском
якорей; низкий bridge rate — уже топологией, температурой, телепортом и типами
допустимых переходов.

\section{Hitting time и доступность далёких зон}

Пусть $T_B$ — время первого попадания в узел $B$. Тогда $\E_A[T_B]$ — \emph{hitting
time} из узла $A$ в $B$. Для памяти это не абстракция: если интересный, но редкий
контрпример имеет слишком большой hitting time относительно типичных режимов сна
или фоновой разведки, то practically он \emph{не существует} для системы.

Отсюда инженерный вывод:
\begin{principle}[Редкий узел без пути — это не «редкое знание», а потерянное знание]
Если значимые узлы достижимы только по длинным и узким маршрутам, память не просто
медленно к ним приходит. Она систематически предпочитает не приходить вовсе, а
значит будет стабильно недокорректированной.
\end{principle}

\section{Спектральная щель и смешивание}

Если $\widetilde P$ эргодична, у неё есть стационарное распределение $\pi$. Скорость
смешивания управляется вторым по модулю собственным значением $\lambda_2$:
\[
\text{чем меньше } |\lambda_2|,\ \text{тем быстрее забывается старт}.
\]
Для памяти это двусмысленное благо.

\begin{itemize}
\item Быстрое смешивание полезно для покрытия и scrubbing.
\item Медленное смешивание полезно для удержания темы и целенаправленного retrieval.
\end{itemize}

Отсюда видно, почему одна и та же политика не может одинаково хорошо служить и
ответу, и сну. Retrieval хочет меньше забывать старт; сон, наоборот, должен уметь
выползать из локальной ямы.

\section{Проводимость и bottleneck starvation}

Пусть $S\subset V$. Проводимость
\[
\Phi(S)=\frac{\sum_{u\in S,\ v\notin S}\pi_u \widetilde P_{uv}}
{\min(\pi(S),\pi(V\setminus S))}
\]
измеряет, насколько легко поток выходит из области $S$. Малые значения $\Phi$
означают bottleneck: память часто крутится внутри района, но плохо ходит наружу.

В прикладном смысле это ведёт к двум симптомам:
\begin{enumerate}
\item cross-cluster fusion почти не случается;
\item контрсвидетельства, сидящие за узким перешейком, плохо добираются до ревизии.
\end{enumerate}

\section{Новая инженерная добавка: два бюджета контекст-пакера}
\label{sec:dualbudget}

Ни в исходном roadmap, ни в базовом дубле это не было выделено достаточно явно,
поэтому здесь вводится как \emph{расширение}, а не как ретушь старого текста.

Пусть retrieval уже вернул множество кандидатов $C$. Обычная ошибка — отдать в
контекст всё, что набрало высокий балл, в одном общем списке. Но для ответа нужны
по меньшей мере два разных ресурса:
\[
B = B_{\mathrm{evidence}} + B_{\mathrm{glue}}.
\]

\begin{definition}[Evidence budget]
Бюджет на узлы-свидетельства: наблюдения, факты, контрфакты, provenance-носители,
то есть то, что действительно несёт доказательную нагрузку.
\end{definition}

\begin{definition}[Glue budget]
Бюджет на связующий материал: переходные концепты, краткие обобщения, локальные
мосты, без которых набор свидетельств останется несвязным.
\end{definition}

Если не разделять эти бюджеты, происходит одно из двух. Либо всё съедает
свидетельственный материал, и ответ становится точным, но несобранным; либо всё
забивается «красивыми мостами», а фактическое основание усыхает. Разделение
бюджетов делает context packer не просто сортировщиком, а маленьким архитектором
локального доказательства.

\section{Что именно надо мерить на стенде}

Из этой главы естественно вытекает retrieval-аудит как кортеж
\[
\mathcal{R}
=
\bigl(
\text{anchor recall},
\text{bridge rate},
\text{contradiction hit rate},
\text{provenance coverage}
\bigr).
\]
Последние два компонента читаются так:

\begin{definition}[Contradiction hit rate]
Доля случаев, в которых при наличии в релевантной окрестности важного контрсвидетельства
retrieval действительно поднял его в рабочий контекст.
\end{definition}

\begin{definition}[Provenance coverage]
Доля включённых в контекст утверждений, для которых рядом лежит хотя бы минимальный
support-skeleton или иной объяснимый след происхождения.
\end{definition}

\begin{workedbox}[Почему одного recall недостаточно]
Пусть ответ на вопрос требует узлов $A$ и $B$, а их связывает мост $M$. Если якорь
нашёлся в районе $A$, а обход принёс $A$ и ещё пять соседей, но не принёс $M$ и $B$,
формально память «что-то помнит», practically — не отвечает. Recall на якорях может
быть высоким при полностью провальном bridge rate.
\end{workedbox}

\section{Вывод}

Математика retrieval нужна не ради строгости как украшения. Она нужна, чтобы не
перепутать четыре разных класса проблем:
\begin{enumerate}
\item ошибка выбора якоря;
\item ошибка локального разворачивания;
\item топологическая недоступность;
\item ошибка упаковки уже найденного материала.
\end{enumerate}

Без такого разделения система с плохим ответом выглядит просто как «что-то не так с
памятью». С ним становится видно, \emph{какой именно слой памяти провалился}.
